\section{Black Hole Interior as a Non-Equilibrium Phase}
\label{sec:black_hole}

\subsection{Revisiting the Meaning of the Horizon}

Section~\ref{sec:schwarzschild} showed that the Schwarzschild radius $r_s = 2GM$
appears as a critical point of the modified distance structure.
At this radius, propagation cost diverges and the exterior solution ceases to apply.
\begin{equation}
  \boxed{\text{The horizon is a phase boundary.}}
\end{equation}
This interpretation differs from the conventional geometric breakdown.
Instead, the horizon marks a transition between two regimes of closure dynamics.

\subsection{Exterior vs.\ Interior}

Outside the horizon, the classical phase holds:
\begin{equation}
  \sigma \approx 0, \qquad \nabla \cdot J \approx 0, \qquad \partial_t C \approx 0.
\end{equation}
In this regime, closure structure is stable, conservation laws hold,
and the metric is well-defined.
\begin{equation}
  \boxed{\text{Exterior = classical closure phase.}}
\end{equation}
Crossing the horizon changes these conditions.

\subsection{Reignition of $\sigma$}

Inside the horizon, the generation term becomes nonzero:
\begin{equation}
  \boxed{\sigma > 0.}
\end{equation}
Log generation resumes and the fixed-point structure breaks down:
$\partial_t C \neq 0$.
This implies that geometry is no longer static.
\begin{equation}
  \boxed{\text{The generative time structure becomes dynamically active again.}}
\end{equation}
The interior therefore corresponds to a non-equilibrium regime of closure.

\subsection{Collapse of the Fixed Point}

In the exterior, $\partial_t C \approx 0$;
in the interior, $\partial_t C \neq 0$.
The metric is no longer determined by a fixed-point condition.
\begin{equation}
  \boxed{\text{The fixed-point structure collapses.}}
\end{equation}
Geometry becomes dynamically generated rather than static.

\subsection{The Singularity as a Closure Limit}

General relativity predicts a singularity at the center.
Within VED,
the same region is interpreted as the limit of gravitational closure.
Using the effective non-closure barrier introduced in Vol.~1,
\begin{equation}
  B(L) = -\kappa_B \log\!\left(1 - \frac{L}{L_{\max}}\right)
  \to \infty
  \qquad (L \to L_{\max}).
\end{equation}
\begin{equation}
  \boxed{\text{Perfect closure is not an ordinary attainable state.}}
\end{equation}
Thus the center is not treated as an ordinary object of infinite density
within the VED descriptive window.
It is read as an asymptotic closure limit.
In the language developed later in Vol.~4,
this limit is a local Differential Horizon of the gravitational
description:
the point where closure attempts to erase the differences required
for the description itself.

\subsection{Interior Dynamics}

With $\sigma > 0$, closure structure is continuously regenerated.
\begin{equation}
  \boxed{\text{The interior is a region of structural regeneration.}}
\end{equation}
Logs are generated, closure reorganizes, and geometry evolves.
The interior is therefore dynamically active rather than a state of collapse.

\subsection{The Inversion}

Conventional interpretation: black hole = endpoint.

VED:
\begin{equation}
  \boxed{\text{Black hole = entry into a new phase.}}
\end{equation}
The exterior is a fixed world; the interior is a world under generation.

\subsection{Correspondence with the Early Universe}

The early universe is characterized by $\sigma > 0$.
The black hole interior also satisfies $\sigma > 0$.
\begin{equation}
  \boxed{
    \text{The early universe and black hole interior are the same phase
    at different scales.}
  }
\end{equation}

\subsection{Conclusion}

\begin{equation}
  \boxed{\text{The black hole interior is not an ordinary singular object.}}
\end{equation}
\begin{equation}
  \boxed{\text{It is a non-equilibrium closure phase.}}
\end{equation}
Closure continues to evolve, and geometry remains dynamically generated.
